\documentclass[conference]{IEEEtran}
\usepackage{amsmath,amssymb}
\usepackage{booktabs}
\usepackage{graphicx}
\usepackage{microtype}
\usepackage{url}
\usepackage[hidelinks]{hyperref}

\title{Product-Aware Explainable Anomaly Screening with Blank-Sky Modulation Diagnostics for XPoSat POLIX Level-2 Observations}

\author{
\IEEEauthorblockN{Author One, Author Two, Author Three, and Project Guide}
\IEEEauthorblockA{Department of [Department Name], [College/University]\\
[City, Country]\\
[corresponding-author email]}
}

\begin{document}
\maketitle

\begin{abstract}
The Polarimeter Instrument in X-rays (POLIX) aboard the X-ray Polarimeter Satellite (XPoSat) produces heterogeneous Level-2 products whose scientific roles differ across exposure, energy-resolved azimuth, source azimuth, light-curve, detector, and roll-resolved measurements. This paper investigates whether those products can support interpretable, label-free observation screening while retaining a scientifically independent physical diagnostic. A product-aware representation was constructed for 25 POLIX Level-2 observations included in the project archive: 10 source observations and 15 blank-sky observations. The deployed 15-feature Matrix-C representation is standardized and analyzed using principal component analysis, KMeans, and Isolation Forest. A model-specific explanation score combines PCA separation, assigned-centroid distance, Isolation Forest feature occlusion, and standardized feature abnormality, then maps influential features to their originating product families. The fixed model flagged four anomaly candidates; three remained flagged across all 100 tested Isolation Forest seeds, whereas the fourth was seed-sensitive. Neutralization tests reproduced five Strong and one Moderate faithfulness verdict over a separate six-case exploratory set. An independent WeightedRoll branch fitted second-harmonic modulation and compared normalized \(Q/U\) coefficients with 13 qualifying blank-sky fits. All ten source raw-modulation values were within the empirical blank-sky range. These results show that archive-relative statistical unusualness and polarization-like modulation evidence are related inspection questions but are not equivalent. The analysis is a small-sample screening study and does not report calibrated polarization degree or official sky polarization angle.
\end{abstract}

\begin{IEEEkeywords}
XPoSat, POLIX, X-ray polarimetry, unsupervised anomaly detection, explainable artificial intelligence, scientific machine learning
\end{IEEEkeywords}

\section{Introduction}

The X-ray Polarimeter Satellite (XPoSat) is an Indian Space Research Organisation mission dedicated to X-ray polarimetry and spectroscopy \cite{isro_xposat_2023,saini2025xposat}. Its Polarimeter Instrument in X-rays (POLIX) is a Thomson-scattering polarimeter whose Level-2 archive exposes several product families rather than a single homogeneous table \cite{polix_handbook_2025,rishin2010thomson,fabiani2018instrumentation}. These products describe complementary properties of an observation: roll-dependent exposure, energy-resolved azimuthal behavior, source azimuth, temporal rate variability, detector balance, and exposure-weighted roll modulation.

This work treats the archive as an applied scientific-computing problem. Each observation is represented by interpretable summary features, screened without anomaly labels, and explained using the fitted model components. A separate physical branch fits the WeightedRoll modulation product and compares source behavior with blank-sky observations in the same archive. This separation is essential: WeightedRoll contains source and background modulation contributions \cite{polix_handbook_2025}, while a multivariate anomaly flag need not indicate polarization-like behavior.

Archive anomaly screening has precedents in astronomy \cite{baron2017weirdest,lochner2021astronomaly}, and PCA, KMeans, and Isolation Forest are established methods \cite{pearson1901pca,macqueen1967kmeans,liu2008isolation}. General local-explanation and anomaly-explanation work supplies relevant concepts \cite{lundberg2017shap,yepmo2022anomaly,yeh2019infidelity}. The reviewed literature, however, did not provide a directly comparable framework combining POLIX Level-2 product-aware representation, model-specific unsupervised explanation, perturbation-based faithfulness checks, and an independent blank-sky-referenced modulation branch. This is a scoped literature finding, not a priority claim.

The research question asks whether heterogeneous POLIX Level-2 products can be represented through scientifically interpretable features and screened using explainable unsupervised learning while retaining a scientifically independent comparison with blank-sky-referenced physical modulation.

The contributions are: (1) a 15-feature representation spanning five product-derived families; (2) complementary geometric, clustering, and isolation-based screening of the 25 archived observations; (3) a four-component local feature score mapped to product provenance; (4) faithfulness, seed, contamination, jackknife, feature-tier, and rank-agreement tests; and (5) an independent WeightedRoll and empirical blank-sky \(Q/U\) diagnostic. The Flask application is an implementation and reproducibility contribution. No result is presented as an astrophysical discovery, official anomaly ground truth, calibrated polarization degree (PD), or official sky polarization angle (PA).

\section{Related Work}

\subsection{XPoSat/POLIX and Stokes Analysis}

Official mission, archive, handbook, and acknowledgment sources define XPoSat, POLIX, and the data-use boundaries \cite{isro_xposat_2023,polix_handbook_2025,pradan_xposat,xposat_ack}. Thomson-scattering polarimetry and its instrumentation context are described in \cite{rishin2010thomson,fabiani2018instrumentation}. Stokes analysis supplies a linear representation using cosine and sine coefficients, but calibrated physical inference requires background and response treatment \cite{kislat2015stokes}.

The handbook describes WeightedRoll as including source and background count-rate modulation and cautions that the released source/blank-sky configuration is insufficient for a polarization measurement \cite{polix_handbook_2025}. We therefore report raw modulation, fitted modulation phase, fit quality, and empirical blank-sky-relative quantities only.

\subsection{Astronomy Anomaly Detection and Explanation}

Unsupervised outlier ranking can prioritize unusual astronomical objects for expert review \cite{baron2017weirdest}. Astronomaly combines anomaly detection with active learning and human feedback \cite{lochner2021astronomaly}. The present sample is smaller, lacks labels, and comprises observations assembled from heterogeneous products; reproducibility and sensitivity are therefore more informative than predictive accuracy.

SHAP is a general explanation framework \cite{lundberg2017shap}, but it is not the method deployed here. The project constructs explanations directly from the fitted PCA, KMeans, and Isolation Forest components. Anomaly-explanation literature distinguishes the explanation of unusual instances from supervised class explanation \cite{yepmo2022anomaly}. Perturbation-based faithfulness concepts motivate checking whether neutralizing influential features reduces model evidence \cite{yeh2019infidelity}.

Multi-view representation learning provides general approaches for combining heterogeneous information sources \cite{hwang2021multiview}. Here, traceable domain-named summaries are retained instead of learning an end-to-end representation because the archive contains only 25 observations.

\section{Data and Problem Formulation}

The data freeze is 27 July 2026. It contains 25 POLIX Level-2 observations included in the project archive: 10 source observations and 15 blank-sky observations. The source set contains three Crab observations and one each for Sco X-1, Her X-1, GX 301-2, Cyg X-1, 4U 1700-37, Cen X-3, and Cas-A SNR. Unique observation identifiers are used as keys.

\begin{table}[t]
\caption{Dataset Composition}
\label{tab:data}
\centering
\begin{tabular}{lrl}
\toprule
Role & Count & Use \\
\midrule
Source & 10 & Screening and physical comparison \\
Blank sky & 15 & Screening and empirical reference \\
Total & 25 & Processed by both branches \\
\bottomrule
\end{tabular}
\end{table}

No trusted anomaly labels exist. ``Anomaly'' therefore means statistically unusual relative to this 25-observation feature distribution. It does not denote a confirmed physical event.

Tier 1A contains exposure and energy-resolved azimuth summaries; Tier 1B contains source-azimuth summaries; Tier 2 contains light-curve and detector diagnostics. WeightedRoll (WR) supplies an exposure-weighted modulation curve containing source and background contributions. WR is excluded from Matrix C and analyzed only in the independent physical branch, preventing circular confirmation.

\section{Product-Aware Feature Engineering}

Four matrices were frozen. Matrix A contains eight exposure and energy-resolved features. Matrix B adds three source-azimuth features. Matrix C adds four temporal and detector features and is the deployed 15-feature input. WR contains six modulation diagnostics and is separate.

\begin{table*}[t]
\caption{Matrix-C Features}
\label{tab:features}
\centering
\small
\begin{tabular}{lll}
\toprule
Family & Feature & Interpretation \\
\midrule
Exposure & \texttt{t1A\_exp\_uniformity\_cv} & Roll-exposure coefficient of variation \\
Exposure & \texttt{t1A\_exp\_max\_to\_min\_roll} & Maximum-to-minimum exposure ratio \\
EnergyRes & \texttt{t1A\_energy\_peak\_channel} & Peak response channel \\
EnergyRes & \texttt{t1A\_energy\_weighted\_mean\_channel} & Weighted channel mean \\
EnergyRes & \texttt{t1A\_energy\_weighted\_std\_channel} & Weighted channel spread \\
EnergyRes & \texttt{t1A\_energy\_high\_channel\_fraction} & High-channel fraction \\
EnergyRes & \texttt{t1A\_energy\_channel\_entropy} & Channel-distribution entropy \\
EnergyRes & \texttt{t1A\_energy\_anode\_balance\_cv} & Across-anode balance variability \\
Source azimuth & \texttt{t1B\_src\_peak\_to\_median\_roll} & Peak-to-median roll ratio \\
Source azimuth & \texttt{t1B\_src\_roll\_entropy} & Roll-distribution entropy \\
Source azimuth & \texttt{t1B\_src\_roll\_smoothness\_norm} & Normalized roll roughness \\
Light curve & \texttt{t2\_lc\_rate\_cv} & Rate coefficient of variation \\
Light curve & \texttt{t2\_lc\_peak\_to\_median\_rate} & Peak-to-median rate \\
Detector LC & \texttt{t2\_det\_lc\_rate\_balance\_cv} & Detector-rate balance variability \\
Detector PHA & \texttt{t2\_det\_pha\_centroid\_spread} & Pulse-height centroid spread \\
\bottomrule
\end{tabular}
\end{table*}

All matrices have 25 rows and no missing values. Although the reviewed report mentioned median imputation, neither the saved artifact nor deployed code contains an imputer. No imputation claim is made. Matrix C was selected by the completed project; the ablation below shows that it is not universally optimal.

\section{Explainable Unsupervised Screening}

\subsection{Model}

The saved \texttt{StandardScaler} transforms each 15-vector \(\mathbf{x}_i\) into \(\mathbf{z}_i\). Two-component PCA explains 38.74\% and 22.51\% of variance. KMeans uses \(k=5\), \texttt{n\_init=20}, and random state 42. Isolation Forest uses 100 trees, contamination 0.16, and random state 42. The Isolation Forest score is not a probability, and contamination is an assumed screening fraction rather than a learned prevalence.

\subsection{Four-Component Attribution}

For feature \(j\), the PCA contribution is
\begin{equation}
p_{ij}=\sum_{r=1}^{2}|z_{ij}w_{rj}|\lambda_r ,
\end{equation}
where \(w_{rj}\) and \(\lambda_r\) are the loading and explained-variance ratio. The KMeans contribution is
\begin{equation}
k_{ij}=(z_{ij}-c_{g(i)j})^2 .
\end{equation}
The Isolation Forest occlusion term is
\begin{equation}
o_{ij}=s(\mathbf{z}_i)-s(\mathbf{z}_i^{(j\leftarrow0)}),
\end{equation}
and standardized abnormality is \(a_{ij}=|z_{ij}|\). With \(N(\cdot)\) denoting within-observation maximum-absolute normalization,
\begin{equation}
e_{ij}=N(p_{ij})+N(k_{ij})+N(\max(o_{ij},0))+N(a_{ij}).
\end{equation}
The top five features are returned and mapped to product families. This heuristic attribution is specific to the fitted stack; it is not SHAP and has no causal interpretation. Explanation sentences are deterministic Python templates, not large-language-model output.

\subsection{Faithfulness}

The exploratory faithfulness set contains six candidates and is distinct from the deployed four. Its highest-ranked standardized features are neutralized to zero. A Strong verdict requires reductions in both PCA distance and Isolation Forest anomaly score; Moderate records a partial response.

\section{Independent Physical Modulation Diagnostics}

WeightedRoll is fitted by inverse-variance weighted least squares:
\begin{equation}
y(\phi)=C+Q\cos(2\phi)+U\sin(2\phi).
\end{equation}
Raw modulation and fitted phase are
\begin{equation}
m_{\mathrm{raw}}=\frac{\sqrt{Q^2+U^2}}{C},\qquad
\psi_{\mathrm{fit}}=\tfrac12\operatorname{atan2}(U,Q).
\end{equation}
The coefficients are fit parameters, not official calibrated sky Stokes products; \(\psi_{\mathrm{fit}}\) is not official sky PA. Fits are acceptable for reduced chi-square \(\chi_\nu^2\le2\), caution for \(2<\chi_\nu^2\le5\), and poor otherwise.

Thirteen of 15 blank-sky fits are acceptable and define the normalized empirical reference
\begin{equation}
\bar q_{\mathrm{blank}}=0.0090920,\quad
\bar u_{\mathrm{blank}}=-0.0064782,
\end{equation}
with sample standard deviations 0.0048569 and 0.0040190. The subtraction and standardized vector distance are archive-specific diagnostics, not official background subtraction. PD sensitivity proxies use assumed \(\mu_{100}=0.40,0.42,0.44\) only; these are not calibrated measurements.

\section{Experimental Protocol}

The fixed artifact was replayed against Matrix C, reproducing predictions, scores, PCA coordinates, KMeans labels, and explanations. For Sco X-1, an exact imported-function audit was adopted as the controlling versioned evidence after it agreed with Notebook 10 and the supplementary reproduction. Seed stability refitted Isolation Forest for seeds 0--99 with the other assumptions fixed. Contamination sensitivity used 0.12, 0.16, 0.20, and 0.24. Leave-one-out analysis refitted the scaler and Isolation Forest on each 24-observation subset. Feature-tier ablation applied the same basic settings to A, B, C, and the separate WR diagnostic. Spearman correlation compared PCA distance, centroid distance, and Isolation Forest score. The stored faithfulness computation was reconstructed from primary artifacts.

The audit environment used Python 3.11.0, NumPy 2.4.6, pandas 3.0.3, scikit-learn 1.9.0 \cite{pedregosa2011sklearn}, SciPy 1.17.1, Astropy 8.0.1, and Flask 3.1.3. The original training versions were not independently recovered because the requirements file is unpinned.

\section{Results}

\subsection{Matrix and Ranking Comparison}

\begin{table}[t]
\caption{Feature-Tier Comparison}
\label{tab:ablation}
\centering
\small
\begin{tabular}{lrrrl}
\toprule
Matrix & \(d\) & PCA var. & Silh. & Candidate suffixes \\
\midrule
A & 8 & .7518 & .4133 & 0018,0020,0006,0005 \\
B & 11 & .7570 & .3586 & 0010,0018,0020,0006 \\
C & 15 & .6125 & .3109 & 0010,0018,0003,0006 \\
WR & 6 & .9484 & .4848 & 0023,0003,0004,0007 \\
\bottomrule
\end{tabular}
\end{table}

Matrix A/B/C Isolation Forest rank correlations are 0.913--0.952, but candidate-set Jaccard agreement is 0.333--0.600. The WR row is not used to confirm Matrix-C candidates. PCA distance and Isolation Forest score agree strongly (\(\rho=0.8946\), \(p=1.64\times10^{-9}\)); KMeans centroid distance has weak rank association with PCA (\(\rho=0.1558\), \(p=0.457\)) and Isolation Forest (\(\rho=0.1866\), \(p=0.372\)).

\subsection{Deployed Candidates and Stability}

The saved model reproduced 21 Normal outputs and four anomaly candidates.

\begin{table*}[t]
\caption{Fixed-Model Candidates and Leading Local XAI Evidence}
\label{tab:candidates}
\centering
\small
\begin{tabular}{lll p{7.2cm} r}
\toprule
ID & Role/target & Score & Leading XAI evidence & Seeds \\
\midrule
C24\_0018 & Blank-13 & .615964 & Energy spread & 100 \\
G01\_0006 & Sco X-1 & .593516 & Peak 2.454499; weighted mean 2.207369; entropy 1.813968 & 100 \\
G01\_0003 & Her X-1 & .572308 & LC rate CV & 100 \\
C24\_0010 & Blank-5 & .517611 & Roll smooth. & 29 \\
\bottomrule
\end{tabular}
\end{table*}

The first three were selected under all four contamination settings and in all 24 included jackknife refits. C24\_0010 was selected in three contamination settings, 19/24 included jackknife fits, and 29/100 seeds. C24\_0020, absent from the fixed deployed four, was selected in 68/100 seeds, two contamination settings, and 7/24 included jackknife fits. The stable result is therefore a three-case core plus an uncertain threshold neighborhood. Common-observation jackknife rankings had median Spearman correlation 0.9896 and minimum 0.9687.

\subsection{Faithfulness and Physical Diagnostics}

Faithfulness metrics matched the stored values within \(9\times10^{-16}\): five verdicts were Strong, one Moderate, and none Weak. C24\_0023 was Moderate because PCA distance fell while the Isolation Forest score did not.

All 25 WeightedRoll products were fitted: 19 acceptable, three caution, and three poor. The 13 acceptable blank-sky fits had mean raw modulation 1.147820\%, sample standard deviation 0.565960\%, median 1.492835\%, and range 0.294127--1.778927\%.

\begin{table*}[t]
\caption{Source Modulation and Blank-Sky-Relative Diagnostics}
\label{tab:physical}
\centering
\small
\begin{tabular}{lrrrrl}
\toprule
Target (ID) & Raw modulation (\%) & \(\chi_\nu^2\) & Fit & Vector distance & Status \\
\midrule
Her X-1 (G01\_0003) & \(0.5966\pm0.0167\) & 57.4335 & Poor & 2.3174 & Moderate; low confidence \\
Sco X-1 (G01\_0006) & \(1.1396\pm0.0199\) & 2.0308 & Caution & 0.3005 & Within scatter \\
Crab (P01\_0005) & \(1.6973\pm0.0310\) & 1.0625 & Acceptable & 1.2649 & Within scatter \\
GX 301-2 (G01\_0004) & \(0.7223\pm0.0295\) & 4.3745 & Caution & 1.2514 & Within scatter \\
Cyg X-1 (T24\_0007) & \(0.6811\pm0.0153\) & 5.4080 & Poor & 1.1480 & Within scatter \\
Crab (T24\_0001) & \(1.3955\pm0.0627\) & 0.8146 & Acceptable & 1.1352 & Within scatter \\
Cas-A SNR (C24\_0026) & \(0.6927\pm0.0294\) & 1.7030 & Acceptable & 0.9385 & Within scatter \\
Cen X-3 (G01\_0002) & \(0.9286\pm0.0220\) & 1.8972 & Acceptable & 0.7748 & Within scatter \\
4U 1700-37 (G01\_0005) & \(1.0806\pm0.0581\) & 1.2291 & Acceptable & 0.6305 & Within scatter \\
Crab (T24\_0002) & \(1.2008\pm0.0220\) & 1.3671 & Acceptable & 0.5152 & Within scatter \\
\bottomrule
\end{tabular}
\end{table*}

All ten source scalar raw-modulation z-scores lie between -0.974 and 0.971 and are within the empirical blank-sky range. These are raw diagnostics, not calibrated PD.

\section{Discussion}

The principal finding is that statistical unusualness and polarization-like modulation evidence are not equivalent. Her X-1 is a stable Matrix-C candidate driven by light-curve variability. Its normalized \(Q/U\) vector has the largest source displacement, but its simple sinusoid fit is poor; this is not clean physical confirmation. Sco X-1 is a stable candidate whose accepted deployed XAI order is energy peak channel (2.454499), energy weighted mean channel (2.207369), and energy channel entropy (1.813968). The historical entropy-first narrative was not reproducible from a located versioned artifact and is superseded by the exact imported-function audit. These three features are local model evidence from the energy-resolved product family and do not establish a physical cause. Sco X-1's vector distance is only 0.3005 and its raw modulation is within blank-sky scatter.

Crab P01\_0005 provides a converse example. It has the largest source raw modulation and an acceptable fit, but is absent from the fixed Matrix-C candidates and appears in only three seed runs. Its vector is still within empirical blank-sky scatter; the value is not a polarization detection.

Blank Sky-13 is a stable candidate explained by energy spread, whereas Blank Sky-5 is a seed-sensitive candidate explained by source-roll smoothness. Blank-sky flags show that the model screens product-level unusualness rather than source physics. Instrumental, exposure, background, processing, or physical causes require targeted investigation and are not assigned here.

The two branches support a triage table: an observation can be statistically unusual, physically displaced, both, or neither. Fit quality adds a necessary gate. The Flask application operationalizes this evidence flow through feature extraction, inference, deterministic explanations, physical diagnostics, visualization, and export, but remains a proof-of-concept deployment.

\section{Limitations and Threats to Validity}

The sample is a project-specific archive of 25 observations. The same archive defines scaling, clusters, isolation structure, and retrospective results; no external holdout is available. Contamination 0.16 is an assumption, feature tier changes the thresholded set, and boundary cases are seed-sensitive. Without labels, accuracy, precision, recall, and superiority claims cannot be made.

The engineered summaries may omit event-level, timing, spectral, calibration, or observation-context information. New releases may cause data drift. WeightedRoll contains source and background contributions; blank-sky behavior varies; and the empirical reference uses only 13 acceptable archive fits. No official source-specific background subtraction, flight-calibrated \(\mu_{100}\), or verified phase-to-sky conversion was available. No domain expert has adjudicated candidates or supplied causal follow-up.

\section{Conclusion}

A reproducible product-aware workflow was demonstrated for the 25 POLIX Level-2 observations included in the project archive. Matrix C supported complementary unsupervised screening; a four-component attribution related observations to features and product families; and perturbation tests reproduced faithfulness verdicts. Stability checks identified a robust three-candidate core while exposing uncertainty at the fixed-model boundary.

The independent WeightedRoll branch fitted raw second-harmonic modulation and compared normalized \(Q/U\) behavior with 13 qualifying blank-sky observations. All ten source raw-modulation values were within the empirical blank-sky range. The comparison shows that archive-relative statistical unusualness does not automatically imply polarization-like modulation, and modulation amplitude does not automatically imply a multivariate anomaly.

The paper does not claim discovery, official anomaly ground truth, official background subtraction, calibrated PD, or official sky PA. Future work should test new POLIX releases, freeze the software environment, use official calibration and background products when available, evaluate prospective drift, and obtain domain-expert review.

\section*{Acknowledgment}

This publication uses the data from the XPoSat mission of the Indian Space Research Organisation (ISRO), archived at the Indian Space Science Data Centre (ISSDC) \cite{xposat_ack}.

The authors also acknowledge [college/project guide/laboratory/funding details to be supplied and approved].

\bibliographystyle{IEEEtran}
\bibliography{08_references}

\end{document}
